\section{Formal Methodology and Causal Engine}
\label{sec:methodology}

\subsection{The Spatio-Temporal Graph}

Let $G = (V, E, W)$ be an undirected weighted graph in which each vertex
$i \in V$, $|V| = N$, is a neighbourhood unit (a census tract in the empirical
setting, a lattice cell in the synthetic setting) and each edge
$(i,j) \in E$ encodes spatial adjacency. We write $\mathcal{N}(i) =
\{ j : (i,j) \in E \}$ for the neighbour set of $i$ and $d_i = |\mathcal{N}(i)|$
for its degree. The weight matrix $W \in \mathbb{R}^{N \times N}$ is the
row-normalised adjacency,
\begin{equation}
  W_{ij} =
  \begin{cases}
    a_{ij} \big/ \sum_{k} a_{ik}, & (i,j) \in E,\\[2pt]
    0, & \text{otherwise,}
  \end{cases}
  \label{eq:weights}
\end{equation}
with $a_{ij}$ the raw adjacency entry, so that $\sum_{j} W_{ij} = 1$ for every
non-isolated $i$. Row normalisation makes the interference term an average
rather than a sum, which keeps its scale invariant to degree and prevents
high-degree units from dominating the objective.

Each unit carries a feature trajectory $X_{i,t} \in \mathbb{R}^{F}$ over
discrete steps $t = 0, \dots, T$. In this work $F = 6$: land value, rent,
population, median income, housing units and accessibility. A policy assignment
$P_{i,t} \in \{0,1\}$ indicates whether unit $i$ is subject to the intervention
at step $t$. We write $X_{i,t-W_h+1:t}$ for the length-$W_h$ history window
ending at $t$.

\subsection{Interference and the Exposure Mapping}

Under SUTVA, the potential outcome of unit $i$ would be written
$Y_i(P_i)$, depending on $i$'s assignment alone. Spatial spillover invalidates
this. We adopt instead the partial interference
assumption~\cite{hudgens2008,ogburn2014}: unit $i$'s outcome depends on its own
assignment and on its neighbours' assignments only through a low-dimensional
exposure summary. We use the degree-normalised exposure
\begin{equation}
  e_i = \sum_{j \in \mathcal{N}(i)} W_{ij} P_{j,t}
      = \frac{1}{d_i} \sum_{j \in \mathcal{N}(i)} P_{j,t},
  \label{eq:exposure}
\end{equation}
so that $e_i \in [0,1]$ measures the treated fraction of $i$'s neighbourhood.
The potential outcome becomes $Y_i(P_i, e_i)$, and the estimands of interest
separate into a direct effect holding exposure fixed and an interference effect
holding own-treatment fixed~\cite{forastiere2021}.

\subsection{Causal Effect Decomposition}

The \CGCEN{} forward pass realises this separation. Let $h_i \in
\mathbb{R}^{H}$ be a node-local temporal encoding of the history window,
produced by a shared feature encoder followed by a gated recurrent
unit~\cite{cho2014gru}:
\begin{align}
  z_{i,k} &= \mathrm{ReLU}\!\left( \mathbf{W}_{\mathrm{enc}} \tilde{X}_{i,k}
             + b_{\mathrm{enc}} \right),
             \quad k = t - W_h + 1, \dots, t, \label{eq:encoder}\\
  h_i     &= \mathrm{GRU}\!\left( z_{i,t-W_h+1}, \dots, z_{i,t} \right),
             \label{eq:gru}
\end{align}
where $\tilde{X}$ denotes the $z$-score standardised features of
Section~\ref{subsec:scaling}. Critically, \eqref{eq:encoder}--\eqref{eq:gru}
contain no message passing: $h_i$ is a function of unit $i$'s own history only.

Writing $s_i = [\, h_i \,\Vert\, P_{i,t} \,]$ for the state-treatment
concatenation, the prediction is
\begin{equation}
  \hat{Y}_{i,t+1}
  = \underbrace{X_{i,t} + \mathrm{MLP}_{\Phi}(s_i)}_{\Phi(X_{i,t}, P_{i,t})}
  + \underbrace{\sum_{j \in \mathcal{N}(i)} W_{ij} \,
    \mathrm{MLP}_{\Psi}(s_j)}_{\sum_j W_{ij} \Psi(X_{j,t}, P_{j,t})} .
  \label{eq:decomposition}
\end{equation}
The term $X_{i,t}$ inside $\Phi$ is a residual anchor: the network predicts the
\emph{change} from the last observation rather than the level. This is not a
cosmetic choice. Without it the GRU must reconstruct an absolute level through a
saturating $\tanh$ bottleneck. Ablating the anchor while holding everything else
fixed degrades pooled MAE from $8.92$ to $37.28$, a factor of $4.2$; the
near-linear ``next $\approx$ last'' signal that dominates these series is
destroyed by the recurrent nonlinearity. Because
$X_{i,t}$ is a function of unit $i$'s own state, the anchor is properly part of
the direct term and does not contaminate the interference channel.

The identification of the split is architectural. $\mathrm{MLP}_{\Phi}$ reads
only $s_i$, which derives from the message-passing-free encoding
\eqref{eq:gru}; it therefore has no path by which neighbour state could enter.
$\mathrm{MLP}_{\Psi}$ is evaluated at the neighbours and aggregated with $W$;
it is the only route by which neighbour information reaches the prediction.

\subsection{Recovering ITE and STE}

Effects are recovered by contrasting each branch against its own untreated
counterfactual. Let $s_i^{0} = [\, h_i \,\Vert\, 0 \,]$. Then
\begin{align}
  \tau_i &= \mathrm{MLP}_{\Phi}(s_i) - \mathrm{MLP}_{\Phi}(s_i^{0}),
  \label{eq:ite}\\
  \eta_i &= \sum_{j \in \mathcal{N}(i)} W_{ij}
            \left( \mathrm{MLP}_{\Psi}(s_j) - \mathrm{MLP}_{\Psi}(s_j^{0}) \right),
  \label{eq:ste}
\end{align}
where $\tau_i$ is the estimated ITE and $\eta_i$ the estimated STE. Two
identities follow immediately and serve as correctness checks in our test
suite. First, if $P_{i,t} = 0$ then $s_i = s_i^{0}$ and $\tau_i = 0$ exactly:
an untreated unit has no direct effect. Second, if $E = \varnothing$ then the
aggregation in \eqref{eq:ste} is empty and $\eta_i = 0$ exactly: with no graph
there is no interference. The additive identity
$\hat{Y}_{i,t+1} = \Phi_i + \sum_j W_{ij}\Psi_j$ holds by construction to
floating-point tolerance.

Algorithm~\ref{alg:causal_stgnn} states the full forward pass together with the
counterfactual contrasts.

\begin{algorithm}[t]
\caption{Causal ST-GNN forward pass with direct/spillover decomposition (CG-CEN)}
\label{alg:causal_stgnn}
\begin{algorithmic}[1]
\REQUIRE Node history $X_{i,t-W+1:t} \in \mathbb{R}^{W \times F}$ for all $i \in \mathcal{V}$;
         policy assignment $P_{i,t} \in \{0,1\}$; adjacency $\mathcal{N}(\cdot)$ with
         row-normalised weights $W_{ij}$
\ENSURE  Prediction $\hat{Y}_{i,t+1}$, individual treatment effect $\tau_i$,
         spatial interference effect $\eta_i$
\STATE $\mu, \sigma \gets$ \textsc{FitScaler}$(X_{\text{train}})$ \COMMENT{training split only}
\STATE $Z_{i,\cdot} \gets (X_{i,\cdot} - \mu) / \sigma$ for all $i$
\FOR{$k = t-W+1$ \TO $t$}
    \STATE $e_{i,k} \gets \mathrm{ReLU}(\mathbf{W}_{\text{enc}} Z_{i,k} + b_{\text{enc}})$
\ENDFOR
\STATE $h_i \gets \mathrm{GRU}\big(e_{i,t-W+1}, \dots, e_{i,t}\big)$ \COMMENT{node-local; no message passing}
\STATE $s_i \gets [\,h_i \,\Vert\, P_{i,t}\,]$
\STATE $\Phi_i \gets \mathrm{MLP}_{\Phi}(s_i) + X_{i,t}$ \COMMENT{direct term, residual anchor}
\STATE $\Psi_j \gets \mathrm{MLP}_{\Psi}(s_j)$ for all $j$
\STATE $\Omega_i \gets \sum_{j \in \mathcal{N}(i)} W_{ij} \, \Psi_j$ \COMMENT{spatial interference}
\STATE $\hat{Y}_{i,t+1} \gets \Phi_i + \Omega_i$
\STATE $s_i^{0} \gets [\,h_i \,\Vert\, 0\,]$ \COMMENT{counterfactual: no intervention}
\STATE $\tau_i \gets \mathrm{MLP}_{\Phi}(s_i) - \mathrm{MLP}_{\Phi}(s_i^{0})$ \COMMENT{ITE}
\STATE $\eta_i \gets \sum_{j \in \mathcal{N}(i)} W_{ij}\big(\mathrm{MLP}_{\Psi}(s_j)
        - \mathrm{MLP}_{\Psi}(s_j^{0})\big)$ \COMMENT{STE}
\RETURN $\hat{Y}_{i,t+1}, \tau_i, \eta_i$
\end{algorithmic}
\end{algorithm}


\subsection{Feature Scaling}
\label{subsec:scaling}

The panel spans several orders of magnitude---median income is $O(10^{4})$
while accessibility is $O(10^{0})$---so an unscaled squared-error objective is
dominated by the large-magnitude columns and the gradient signal for the
remainder is effectively lost. We apply per-feature $z$-score standardisation,
$\tilde{X}_{\cdot,\cdot,f} = (X_{\cdot,\cdot,f} - \mu_f)/\sigma_f$, with
$\mu_f$ and $\sigma_f$ estimated \emph{on the training split only} and reused
unchanged for the test split, so that no distributional information leaks
backwards across the temporal boundary. Constant features are assigned
$\sigma_f = 1$ to avoid division by zero. All reported metrics are computed
after inverting the transform, so errors are in the original units and remain
comparable across configurations.

\subsection{Baseline Models}

We compare against three baselines spanning the relevant hypothesis space.
The \emph{Temporal MLP} flattens the history window and maps it to the next
step through a two-layer perceptron; it is graph-blind by construction and
isolates the value of the topology. The \emph{Spatial Lag Model} (SLM) is a
learned analogue of the classical spatial autoregressive specification of
spatial econometrics~\cite{anselin1988,lesage2009}, predicting
\begin{equation}
  \hat{Y}_{i,t+1} = X_{i,t} + \beta^{\top} \mathrm{vec}(X_{i,t-W_h+1:t})
  + \rho \odot \sum_{j \in \mathcal{N}(i)} W_{ij} X_{j,t},
  \label{eq:slm}
\end{equation}
with a learned spatial coefficient $\rho$ and a residual anchor matching
\eqref{eq:decomposition} so that the comparison isolates functional form rather
than parameterisation. The \emph{Naive Persistence} model predicts
$\hat{Y}_{i,t+1} = X_{i,t}$; it is parameter-free and never trained. It is the
floor every learned model must clear before any claim about spillover structure
is meaningful, and including it is what revealed, in an earlier iteration of
this work, that two learned models were both losing to it.

\subsection{Policy Indices}

\subsubsection{Housing $+$ Transportation Overburden Index}

Housing cost alone understates the burden borne by households in
poorly-connected neighbourhoods, because low rents are frequently purchased
with high commuting costs. Following the combined housing-and-transportation
affordability convention, we define
\begin{equation}
  \HT_{i,t}
  = \frac{\mathrm{Rent}_{i,t} + \mathrm{TransitCost}_{i,t}}
         {\mathrm{MedianIncome}_{i,t}} ,
  \label{eq:ht}
\end{equation}
where transit cost falls with accessibility,
\begin{equation}
  \mathrm{TransitCost}_{i,t} =
  \frac{c_0}{\max(\mathrm{Accessibility}_{i,t}, \epsilon)^{\kappa}} ,
  \label{eq:transit}
\end{equation}
with base cost $c_0$ and elasticity $\kappa$. This is precisely the mechanism a
rent-to-income ratio misses: a transit investment that raises accessibility
lowers $\HT$ even when it raises rent, and whether the net effect is
progressive or regressive is an empirical question the index can answer.

\subsubsection{Level-Aware Displacement Pressure Index}

The Displacement Pressure Index blends adverse movements in rent, low-income
share and accessibility:
\begin{align}
  \DPI_{i,t} =\;& \alpha \left( \frac{\Delta \mathrm{Rent}_{i,t}}
                                    {\mathrm{Rent}_{i,t-1}} \right)
  + \beta \left( 1 - \frac{\mathrm{LowIncomeShare}_{i,t}}
                          {\mathrm{LowIncomeShare}_{i,t-1}} \right) \nonumber\\
  &+ \gamma \left( \Delta \mathrm{Accessibility}_{i,t} \right)
   + \delta \, B_{i,t} .
  \label{eq:dpi}
\end{align}
The first three terms are the growth-rate formulation. The fourth is a
correction we introduce. A growth-rate index is, by construction, blind to
sustained level differences: a policy that permanently raises rents by six per
cent changes the growth rate only in the single period in which the shift
occurs, after which the index returns to its baseline path. In our experiments
this rendered every policy contrast on the order of $10^{-4}$---numerically
present but substantively meaningless. We therefore add a cumulative
rent-burden level term measured against each unit's own initial burden,
\begin{equation}
  B_{i,t} = \frac{\HT_{i,t}}{\max(\HT_{i,0}, \epsilon)} - 1 ,
  \label{eq:burden}
\end{equation}
which is zero at $t = 0$ by construction and retains a permanent offset under a
permanent shift. Setting $\delta = 0$ recovers the pure growth-rate
specification exactly, a property we pin with a regression test so that the
original formulation remains verifiable.

Both indices are constructed comparative indicators over a controlled panel.
Neither is a calibrated forecast of realised displacement, and we use them only
for ranking scenarios against a common baseline.
